% 人工对照底本逐节点转写；独立TikZ重绘。

\begin{tikzpicture}[x=1cm,y=-1cm,every node/.style={anchor=west,align=left,inner sep=2pt,font=\fontsize{9}{11.7}\selectfont},edge/.style={draw=black!65,line width=.45pt}]
\node[text width=0.753cm,draw=black!65] (n1) at (0.000,4.607) {客观\\要件};
\node[text width=1.350cm,draw=black!65] (n2) at (1.420,0.477) {没有危害\\行为};
\node[text width=2.809cm,draw=black!65] (n3) at (3.190,0.477) {不能犯（1号案例）};
\node[text width=1.069cm,draw=black!65] (n4) at (1.420,5.199) {有危害\\行为};
\node[text width=0.753cm,draw=black!65] (n5) at (3.190,5.199) {主观\\要件};
\node[text width=5.972cm,draw=black!65] (n6) at (4.710,1.459) {缺乏犯罪故意型事实认识错误（2号案例）};
\node[text width=2.334cm,draw=black!65] (n7) at (4.710,5.589) {具有犯罪故意型\\事实认识错误};
\node[text width=1.702cm,draw=black!65] (n8) at (7.480,4.997) {具体的错误};
\node[text width=1.860cm,draw=black!65] (n9) at (9.800,2.441) {对象错误\\（3号案例）};
\node[text width=1.702cm,draw=black!65] (n10) at (12.270,2.441) {无观点争议};
\node[text width=1.385cm,draw=black!65] (n11) at (9.800,3.813) {打击错误};
\node[text width=1.702cm,draw=black!65] (n12) at (12.270,3.423) {具体符合说};
\node[text width=1.702cm,draw=black!65] (n13) at (12.270,4.203) {法定符合说};
\node[text width=1.385cm,draw=black!65] (n14) at (9.800,6.369) {因果关系\\错误};
\node[text width=3.283cm,draw=black!65] (n15) at (12.270,5.185) {狭义的因果关系错误。\\重点是因果关系};
\node[text width=2.651cm,draw=black!65] (n16) at (12.270,6.369) {结果的推迟发生。\\重点是因果关系};
\node[text width=2.651cm,draw=black!65] (n17) at (12.270,7.553) {结果的提前实现。\\重点是着手};
\node[text width=1.860cm,draw=black!65] (n18) at (7.480,8.737) {抽象的错误\\（4号案例）};
\node[text width=4.232cm,draw=black!65] (n19) at (9.800,8.737) {分别分析两个罪，想象竞合，\\注意包容评价关系};
\draw[edge] (n2.east)--(3.040,0.477)--(3.040,0.477)--(n3.west);
\draw[edge] (n1.east)--(1.270,4.607)--(1.270,0.477)--(n2.west);
\draw[edge] (n5.east)--(4.560,5.199)--(4.560,1.459)--(n6.west);
\draw[edge] (n9.east)--(12.120,2.441)--(12.120,2.441)--(n10.west);
\draw[edge] (n8.east)--(9.650,4.997)--(9.650,2.441)--(n9.west);
\draw[edge] (n11.east)--(12.120,3.813)--(12.120,3.423)--(n12.west);
\draw[edge] (n11.east)--(12.120,3.813)--(12.120,4.203)--(n13.west);
\draw[edge] (n8.east)--(9.650,4.997)--(9.650,3.813)--(n11.west);
\draw[edge] (n14.east)--(12.120,6.369)--(12.120,5.185)--(n15.west);
\draw[edge] (n14.east)--(12.120,6.369)--(12.120,6.369)--(n16.west);
\draw[edge] (n14.east)--(12.120,6.369)--(12.120,7.553)--(n17.west);
\draw[edge] (n8.east)--(9.650,4.997)--(9.650,6.369)--(n14.west);
\draw[edge] (n7.east)--(7.330,5.589)--(7.330,4.997)--(n8.west);
\draw[edge] (n18.east)--(9.650,8.737)--(9.650,8.737)--(n19.west);
\draw[edge] (n7.east)--(7.330,5.589)--(7.330,8.737)--(n18.west);
\draw[edge] (n5.east)--(4.560,5.199)--(4.560,5.589)--(n7.west);
\draw[edge,->] (n4.east)--(3.040,5.199)--(3.040,5.199)--(n5.west);
\draw[edge] (n1.east)--(1.270,4.607)--(1.270,5.199)--(n4.west);
\end{tikzpicture}

